% =============================================================================
% Discussion 改写 — 自然融入审稿人关切，不显回复痕迹
% 操作: 全文替换原文 Section 4 (Discussion)
% =============================================================================

\section{Discussion}

This study evaluated \dentalclaw{} as a natural-language-driven workflow platform
for dental imaging AI research. The results suggest that \dentalclaw{} can convert
user research instructions into executable workflows covering dataset preparation,
model training, inference, evaluation, and structured reporting. Beyond confirming
end-to-end executability, we further assessed the platform's intent-to-workflow
decision accuracy through a 6-dimension weighted evaluation across 11~representative
intents, achieving a weighted score of 0.977 (Section~3.4). Importantly, the
contribution of \dentalclaw{} should be interpreted as workflow infrastructure
rather than as a new segmentation model: its value lies in organising fragmented,
operator-dependent steps into a standardised, reviewable, and auditable process for
dental AI experimentation.

\paragraph{Relationship to existing frameworks.}
\dentalclaw{} is best understood as an \textit{orchestration layer} that sits above
specialised deep learning frameworks, rather than as a competitor to them. MONAI
provides a comprehensive library of medical imaging components (transforms, network
architectures, loss functions) for building custom training pipelines; nnU-Net
automates the configuration of segmentation models from dataset statistics. Both
are powerful tools for the \textit{how} of model development. \dentalclaw{} addresses
a different, complementary problem: the \textit{what} and \textit{whether}---given a
research question expressed in natural language, which of several possible workflows
should be executed, with which dataset, under what configuration, and with what
quality-control checks? The platform does not replace MONAI or nnU-Net; it delegates
model execution to them (or to registered skill scripts) after making an informed
selection. This separation of concerns---decision from execution---is the
architectural premise that distinguishes \dentalclaw{} from existing workflow
managers such as Snakemake or Nextflow, which require users to pre-specify workflow
logic in a domain-specific language.

\paragraph{Limitations and failure modes.}
The agent's decision can degrade in two predictable ways, both of which are handled
by the platform's design. When both the dataset and the task type are unrecognised
from the user's input, the platform returns an out-of-scope rejection---this is a
safe failure because no workflow is executed. When only one of these fields is
recognised, the platform asks for clarification rather than guessing the missing
information. A more subtle failure occurs when the agent proposes an external
solution with an incorrect or non-existent entrypoint (e.g., a hallucinated script
path). This risk is mitigated by design: external proposals carry an
\texttt{external\_suggestion} status and are \textit{never executed automatically};
they require an explicit user flag. In our evaluation, the two boundary intents that
triggered external proposals (training and detection) produced actionable suggestions
with confidence scores of 0.85 and 0.70, respectively, and neither was executed
without human review. A formal user study with dental researchers would be necessary
to quantify how often such proposals prove useful in practice.

\paragraph{Intent-to-workflow validation.}
A central concern for any agent-driven system is whether the decision pipeline
itself is reliable. The 6-dimension evaluation in Section~3.4 provides quantitative
evidence that the platform correctly distinguishes executable, ambiguous, and
out-of-scope requests in the tested scenarios. The QC~Blocking and Ambiguity~Handling
dimensions both achieved perfect scores (1.000), confirming that the platform never
silently executed an incorrect workflow. The Planning dimension scored 0.955, with
the only deviation arising from boundary cases where no registered method existed
and the agent appropriately proposed external alternatives. These results, while
encouraging, are limited to the 11~intents tested and to the specific agent
configuration (DeepSeek~V4~Pro with OpenClaw). Broader benchmarking across more
diverse intents, agent backends, and language inputs is needed before the decision
pipeline can be considered robust.

\paragraph{Scope of the present validation.}
The present study was designed as a platform validation rather than a clinical
evaluation. The structured reports generated by \dentalclaw{} are intended as
research-audit artefacts (recording what was done, with which data, under what
configuration, and with what outcome), not as diagnostic or treatment-planning
tools. Consistent with this scope, the reported metrics (Dice, workflow completion
rates, decision accuracy scores) reflect technical performance rather than clinical
efficacy. The two representative workflows (2D~panoramic and 3D~CBCT segmentation)
were chosen to demonstrate the platform's ability to handle methodologically distinct
imaging scenarios under a shared infrastructure, not to establish state-of-the-art
segmentation results. The reduction in manual coordination burden compared with
manually assembled scripts is reported as a descriptive observation; a controlled
usability study with clinical researchers would be required to quantify
accessibility improvements.

\dentalclaw{} completed end-to-end workflows in two methodologically different
scenarios: 2D~panoramic radiograph segmentation and 3D~CBCT segmentation. These
tasks differ in data dimensionality, file organisation, preprocessing requirements,
and model execution logic. The ability to handle both workflows under a shared
platform structure suggests that a case-level data representation, combined with
task-specific execution modules and an agent-driven decision layer, can support
standardised AI experimentation across heterogeneous dental imaging modalities.

In summary, the present findings support the feasibility of an agent-driven,
natural-language-programmable platform for dental imaging AI research. The platform's
decision pipeline was validated on a representative set of intents, its limitations
and failure modes were characterised, and its relationship to existing frameworks
was clarified. Further validation through broader task coverage, multi-backend agent
comparisons, and user studies involving clinicians and dental researchers will be
necessary to determine the extent to which such a platform improves accessibility,
efficiency, and reproducibility in real research settings.
